% -----------------------------------------------------------------------------
% Conteudo AODV-EN. Template base: Prof. Wesley Pacheco Calixto (IFG).
% -----------------------------------------------------------------------------
\chapter{REFERENCIAL TEÓRICO}\label{cap:referencial}\label{referencial-teuxf3rico}

Este capítulo apresenta os fundamentos teóricos que embasam o
desenvolvimento da pesquisa. Inicialmente, são discutidos os conceitos
de Internet das Coisas e de Redes de Sensores Sem Fio, situando o
contexto tecnológico do estudo. Em seguida, são apresentadas as
características do microcontrolador ESP32 e do protocolo ESP-NOW,
utilizados na implementação do artefato. Posteriormente, abordam-se os
conceitos de redes mesh e a classificação dos protocolos de roteamento
ad-hoc. Por fim, é analisado o protocolo AODV, que fundamenta a
adaptação proposta e o desenvolvimento do protocolo AODV-EN.

\section{Internet das Coisas e Redes de Sensores Sem
Fio}\label{internet-das-coisas-e-redes-de-sensores-sem-fio}

A Internet das Coisas (\emph{Internet of Things} -- IoT) representa um
paradigma tecnológico no qual objetos físicos dotados de capacidade
computacional, sensores e conectividade de rede integram-se ao ambiente
de forma ubíqua, coletando e compartilhando dados sem intervenção humana
direta \cite{priyadarshi2025}. Essa visão, antecipada por Mark Weiser em
1991 sob o conceito de Computação Ubíqua \cite{weiser1991}, materializa-se atualmente em
bilhões de dispositivos conectados mundialmente, permeando setores como
saúde, agricultura, indústria, cidades inteligentes e automação
residencial. As Redes de Sensores Sem Fio (\emph{Wireless Sensor
Networks} -- WSNs) constituem uma das principais infraestruturas para
implementação de aplicações IoT. Uma WSN é composta por um conjunto de
nós sensores distribuídos espacialmente, capazes de monitorar fenômenos
físicos ou ambientais, como temperatura, umidade, pressão, movimento ou
luminosidade, e transmitir os dados coletados por meio de enlaces sem fio
até uma estação base ou gateway \cite{chaizeng2021}. Os nós sensores são
tipicamente dispositivos de recursos limitados, com restrições de
processamento, memória, energia e largura de banda de comunicação. O
projeto de WSNs deve considerar um conjunto de requisitos críticos que
influenciam diretamente a viabilidade e eficácia das aplicações. O
Quadro~\ref{quadro:req-wsn} sintetiza os principais requisitos e suas implicações para o
projeto de redes de sensores.

{\def\LTcaptype{quadro}
\begin{longtable}[]{@{}
  >{\centering\arraybackslash}p{(\linewidth - 2\tabcolsep) * \real{0.5000}}
  >{\centering\arraybackslash}p{(\linewidth - 2\tabcolsep) * \real{0.5000}}@{}}
\caption{Requisitos críticos para projeto de Redes de Sensores Sem Fio}\label{quadro:req-wsn}\\
\toprule\noalign{}
\begin{minipage}[b]{\linewidth}\centering
Requisito
\end{minipage} & \begin{minipage}[b]{\linewidth}\centering
Descrição e Implicações
\end{minipage} \\
\midrule\noalign{}
\endhead
\bottomrule\noalign{}
\endlastfoot
Baixo custo & Implantações em larga escala requerem dispositivos
economicamente acessíveis. O custo por nó deve permitir a instalação de
dezenas a centenas de dispositivos sem inviabilizar o projeto. \\
Baixo consumo energético & Muitas aplicações exigem operação por meses
ou anos com baterias de capacidade limitada. O consumo deve ser
minimizado em todos os subsistemas: processamento, comunicação e
sensoriamento. \\
Confiabilidade & A rede deve garantir a entrega dos dados coletados
mesmo em condições adversas, como interferência eletromagnética,
obstáculos físicos ou falha de nós individuais. \\
Escalabilidade & A arquitetura deve suportar o crescimento do número de
nós sem degradação significativa de desempenho ou necessidade de
reconfiguração manual. \\
Auto-organização & Os nós devem ser capazes de estabelecer e manter a
topologia da rede de forma autônoma, adaptando-se a mudanças no ambiente
ou na composição da rede. \\
\end{longtable}
}

Fonte: Adaptado de \citeonline{priyadarshi2025} e \citeonline{chaizeng2021}.

A comunicação sem fio é o componente de maior consumo energético em um
nó sensor típico, podendo representar até 70\% do consumo total do
dispositivo \cite{priyadarshi2025}. Dessa forma, a escolha do protocolo
de comunicação e a estratégia de roteamento adotada impactam diretamente
a vida útil da rede. Protocolos que minimizam o número de transmissões,
utilizam modos de baixo consumo durante períodos de inatividade e
otimizam o tamanho das mensagens de controle contribuem
significativamente para a extensão da autonomia dos dispositivos.

\section{O Microcontrolador ESP32}\label{o-microcontrolador-esp32}

O ESP32 é um \emph{System-on-Chip} (SoC) desenvolvido pela Espressif
Systems, projetado para aplicações de Internet das Coisas que demandam
conectividade sem fio integrada, desempenho computacional elevado e
eficiência energética. Desde seu lançamento em 2016, o ESP32 tornou-se
uma das plataformas mais populares para prototipagem e desenvolvimento
de soluções IoT, sendo utilizado tanto em projetos acadêmicos quanto em
produtos comerciais \cite{becker2025}. A arquitetura do ESP32 baseia-se
em um processador dual-core Xtensa LX6, operando em frequências de até
240 MHz, complementado por coprocessador de ultra-baixo consumo (ULP),
memória SRAM de 520 KB e suporte a memória flash externa de até 16 MB. O
SoC integra interfaces Wi-Fi 802.11 b/g/n e Bluetooth 4.2 (incluindo
BLE), além de um amplo conjunto de periféricos como conversores
analógico-digitais (ADC), conversores digital-analógicos (DAC),
interfaces SPI, I²C, I²S, UART, PWM, sensores de toque capacitivo e
sensor de efeito Hall. O Quadro~\ref{quad:esp32-specs} sintetiza as principais especificações técnicas do dispositivo.

O Quadro~\ref{quadro:esp32-specs} resume as principais especificações técnicas do ESP32.

{\def\LTcaptype{quadro}
\begin{longtable}[]{@{}cl@{}}
\caption{Especificações técnicas do microcontrolador ESP32}\label{quadro:esp32-specs}\\
\toprule\noalign{}
Característica & Especificação \\
\midrule\noalign{}
\endhead
\bottomrule\noalign{}
\endlastfoot
Processador & Dual-core Xtensa LX6, até 240 MHz \\
Coprocessador & ULP (Ultra Low Power) para operação em deep sleep \\
Memória SRAM & 520 KB \\
Flash externa & Até 16 MB (tipicamente 4 MB em módulos DevKit) \\
Conectividade Wi-Fi & IEEE 802.11 b/g/n, 2.4 GHz, até 150 Mbps \\
Conectividade Bluetooth & Bluetooth 4.2 BR/EDR e BLE \\
Tensão de operação & 3.0 V -- 3.6 V (típico 3.3 V) \\
GPIOs & 34 pinos programáveis \\
Interfaces & SPI, I²C, I²S, UART, PWM, ADC (12-bit), DAC (8-bit) \\
\end{longtable}
}

Fonte: \citeonline{espressif2024}.

Uma característica diferencial do ESP32 para aplicações IoT é seu
suporte a múltiplos modos de operação com diferentes perfis de consumo
energético. O Quadro~\ref{quadro:esp32-energia} apresenta os modos de energia disponíveis e seus
respectivos consumos típicos, conforme especificado no \emph{datasheet}
do fabricante.

{\def\LTcaptype{quadro}
\begin{longtable}[]{@{}
  >{\raggedright\arraybackslash}p{(\linewidth - 4\tabcolsep) * \real{0.3333}}
  >{\raggedright\arraybackslash}p{(\linewidth - 4\tabcolsep) * \real{0.3333}}
  >{\raggedright\arraybackslash}p{(\linewidth - 4\tabcolsep) * \real{0.3333}}@{}}
\caption{Modos de operação e consumo energético do ESP32}\label{quadro:esp32-energia}\\
\toprule\noalign{}
\begin{minipage}[b]{\linewidth}\raggedright
Modo de Operação
\end{minipage} & \begin{minipage}[b]{\linewidth}\raggedright
Consumo Típico
\end{minipage} & \begin{minipage}[b]{\linewidth}\raggedright
Descrição
\end{minipage} \\
\midrule\noalign{}
\endhead
\bottomrule\noalign{}
\endlastfoot
Active (Wi-Fi TX) & 160 -- 260 mA & Transmissão ativa de dados via
Wi-Fi \\
Active (CPU @ 240 MHz) & 30 -- 68 mA & CPU operando em máxima
frequência, rádio desligado \\
Active (CPU @ 80 MHz) & 20 -- 25 mA & CPU operando em frequência
reduzida, rádio desligado \\
Modem Sleep & 3 -- 20 mA & CPU ativa, rádio Wi-Fi/BT desligado \\
Light Sleep & 0.8 mA & CPU pausada, RAM preservada, despertar rápido
(\textless1 ms) \\
Deep Sleep & 10 -- 150 µA & CPU desligada, apenas RTC ativo, despertar
por timer ou GPIO \\
Hibernation & 5 µA & Apenas RTC timer ativo, memória não preservada \\
\end{longtable}
}

Fonte: \citeonline{espressif2024}.

A disponibilidade de módulos de desenvolvimento (DevKit) a preços
acessíveis, tipicamente entre R\$ 25,00 e R\$ 40,00 no mercado
brasileiro em 2024, democratiza o acesso à tecnologia e viabiliza a construção
de redes de sensores com dezenas de nós a custos compatíveis com
aplicações em países em desenvolvimento.

\section{O Protocolo ESP-NOW}\label{o-protocolo-esp-now}

O ESP-NOW é um protocolo de comunicação sem fio proprietário
desenvolvido pela Espressif Systems, projetado para permitir a troca
direta de mensagens curtas entre dispositivos ESP32 e ESP8266 sem a
necessidade de estabelecimento prévio de conexão Wi-Fi ou associação a
um ponto de acesso. O protocolo opera na camada de enlace do padrão IEEE
802.11, utilizando \emph{action frames} do tipo \emph{vendor-specific}
para transmissão de dados \cite{becker2025}.

\subsection{Características
técnicas}\label{caracteruxedsticas-tuxe9cnicas}

O ESP-NOW opera na frequência de 2.4 GHz, com taxa de transmissão de 1
Mbit/s e largura de banda de canal de 20 MHz. O protocolo implementa um
mecanismo de acesso ao meio baseado em CSMA (\emph{Carrier Sense
Multiple Access}) do tipo \emph{p-persistent}, com \emph{slots}
temporais de aproximadamente 481 µs. O suporte a retransmissões
automáticas permite até 31 tentativas por pacote em caso de falha na
entrega, aumentando a confiabilidade da comunicação \cite{becker2025}.
Cada mensagem ESP-NOW pode transportar até 250 bytes de \emph{payload}
útil, com um \emph{overhead} de cabeçalho de aproximadamente 43 bytes. O
protocolo suporta criptografia opcional utilizando o algoritmo AES-128
no modo CCMP, permitindo comunicação segura quando necessário. A
identificação dos dispositivos é realizada através de seus endereços MAC
de 6 bytes. O Quadro~\ref{quad:espnow-specs} resume as características técnicas do protocolo.

O Quadro~\ref{quadro:espnow-carac} sintetiza as principais características técnicas do ESP-NOW.

{\def\LTcaptype{quadro}
\begin{longtable}[]{@{}
  >{\centering\arraybackslash}p{(\linewidth - 2\tabcolsep) * \real{0.5000}}
  >{\centering\arraybackslash}p{(\linewidth - 2\tabcolsep) * \real{0.5000}}@{}}
\caption{Características técnicas do protocolo ESP-NOW}\label{quadro:espnow-carac}\\
\toprule\noalign{}
\begin{minipage}[b]{\linewidth}\centering
Parâmetro
\end{minipage} & \begin{minipage}[b]{\linewidth}\centering
Valor
\end{minipage} \\
\midrule\noalign{}
\endhead
\bottomrule\noalign{}
\endlastfoot
Padrão base & IEEE 802.11 b/g/n (\emph{action frames} vendor-specific) \\
Frequência de operação & 2.4 GHz (canais 1-14) \\
Taxa de transmissão & 1 Mbit/s \\
Payload máximo & 250 bytes \\
Overhead de cabeçalho & \textasciitilde43 bytes \\
Mecanismo de acesso ao meio & CSMA p-persistent (slot \textasciitilde481
µs) \\
Retransmissões automáticas & Até 31 tentativas por pacote \\
Criptografia & AES-128 CCMP (opcional) \\
Número máximo de peers & 20 peers no total, dos quais até 6 com criptografia habilitada \cite{espressif2024} \\
\end{longtable}
}

Fonte: \citeonline{becker2025}; \citeonline{espressif2024}.

\subsection{Desempenho e alcance}\label{desempenho-e-alcance}

Estudos experimentais recentes fornecem caracterização detalhada do
desempenho do ESP-NOW em diferentes condições. \citeonline{becker2025},
em experimentos de campo conduzidos em 246 posições distintas em
ambientes abertos (fazenda) e florestais, reportaram latência base de
transmissão entre 2.700 e 2.800 µs (aproximadamente 2,8 ms), podendo
alcançar até 60 ms em cenários com múltiplas retransmissões. Em termos
de confiabilidade, foi observado \emph{Packet Delivery Ratio} (PDR)
superior a 99\% em distâncias de até 55 metros em linha de visada, com
alcance máximo de aproximadamente 70 metros antes da perda total de
conectividade. Em ambiente interno, \citeonline{urazayev2023} avaliaram
o desempenho do ESP-NOW considerando a presença de paredes e obstáculos.
Os resultados demonstraram que o ESP-NOW apresenta alcance interno de 15
a 90 metros dependendo das condições, com PDR 15\% superior e consumo
energético 30\% inferior quando comparado a soluções baseadas em Wi-Fi
TCP convencional. Os autores destacam que fatores ambientais como
reflexões e obstáculos afetam significativamente o desempenho, e que
pequenas variações na posição do receptor podem causar grandes
alterações na qualidade do enlace.

\subsection{Limitações para comunicação
multi-hop}\label{limitauxe7uxf5es-para-comunicauxe7uxe3o-multi-hop}

Apesar de suas vantagens, o ESP-NOW apresenta limitações estruturais que
restringem sua aplicação direta em redes de maior escala. A primeira e
mais significativa limitação é a ausência de suporte nativo a roteamento
em múltiplos saltos (\emph{multi-hop}), isto é, ao encaminhamento de
pacotes por nós intermediários até um destino fora do alcance direto. O
protocolo foi concebido exclusivamente para comunicações
ponto-a-ponto ou em topologia estrela, não oferecendo mecanismos para
esse encaminhamento \cite{becker2025}. A segunda limitação refere-se ao número máximo de
\emph{peers} simultâneos: cada dispositivo ESP32 suporta no máximo
20 \emph{peers} registrados na tabela de pares, sendo que apenas 6
desses podem ter criptografia AES-128 habilitada individualmente
\cite{espressif2024}. Essa restrição, imposta por limitações de memória
do hardware, impede a escalabilidade direta da rede e exige estratégias
de gerenciamento dinâmico de \emph{peers} para contorná-la. A terceira limitação relaciona-se à
ausência de mecanismo de \emph{broadcast} nativo eficiente. Embora seja
possível enviar mensagens para o endereço de \emph{broadcast}
(FF:FF:FF:FF:FF:FF), essa abordagem não oferece confirmação de entrega e
apresenta menor confiabilidade, dificultando sua utilização para
disseminação de mensagens de controle em protocolos de descoberta de
rotas.

\section{Redes Mesh e Protocolos de Roteamento
Ad-Hoc}\label{redes-mesh-e-protocolos-de-roteamento-ad-hoc}

\subsection{Conceitos fundamentais de redes
mesh}\label{conceitos-fundamentais-de-redes-mesh}

As Redes Mesh Sem Fio (\emph{Wireless Mesh Networks} -- WMNs) constituem
uma arquitetura de rede descentralizada na qual cada nó pode funcionar
simultaneamente como origem, destino e roteador de dados, estabelecendo
múltiplos caminhos de comunicação entre os participantes
\cite{chaizeng2021}. Diferentemente das topologias centralizadas, como a
estrela, onde todos os nós dependem de um ponto de acesso único, as
redes mesh distribuem a responsabilidade de encaminhamento entre todos
os participantes, eliminando pontos únicos de falha e aumentando a
resiliência do sistema. A comunicação em múltiplos saltos
(\emph{multi-hop}) é uma característica fundamental das redes mesh:
quando o destino está fora do alcance direto de transmissão do nó de
origem, pacotes de dados são encaminhados através de nós intermediários
que atuam como \emph{relays}. Essa abordagem permite estender
significativamente o alcance efetivo da rede sem aumentar a potência de
transmissão individual dos dispositivos, característica particularmente
valiosa para redes de sensores alimentadas por bateria. As WMNs podem
ser classificadas em três categorias principais conforme sua
arquitetura: (i) \emph{Infrastructure WMN}, onde roteadores mesh formam
a infraestrutura de backbone da rede; (ii) \emph{Client WMN}, onde os
próprios dispositivos clientes participam do roteamento; e (iii)
\emph{Hybrid WMN}, que combina características de ambas as abordagens.
Para redes de sensores com ESP32, a arquitetura \emph{Client WMN} é
tipicamente mais adequada, pois todos os nós possuem capacidade
computacional similar e podem contribuir para o roteamento
\cite{chaizeng2021}.

\subsection{Classificação dos protocolos de
roteamento}\label{classificauxe7uxe3o-dos-protocolos-de-roteamento}

Os protocolos de roteamento para redes ad-hoc e mesh podem ser
classificados em três categorias principais conforme sua estratégia de
descoberta e manutenção de rotas: proativos (\emph{table-driven}),
reativos (\emph{on-demand}) e híbridos. Os protocolos proativos, como o
\emph{Optimized Link State Routing} (OLSR) e o
\emph{Destination-Sequenced Distance-Vector} (DSDV), mantêm tabelas de
roteamento constantemente atualizadas através da troca periódica de
informações de topologia entre todos os nós da rede. Essa abordagem
oferece baixa latência na descoberta de rotas --- uma vez que a
informação já está disponível quando necessária --- mas impõe
\emph{overhead} de controle constante, mesmo quando não há tráfego de
dados, o que pode ser problemático para redes de sensores com recursos
energéticos limitados. Os protocolos reativos, como o \emph{Ad-hoc
On-Demand Distance Vector} (AODV) e o \emph{Dynamic Source Routing}
(DSR), descobrem rotas apenas quando há necessidade efetiva de
comunicação. Quando um nó precisa enviar dados para um destino
desconhecido, inicia um processo de descoberta de rota através da
disseminação de mensagens de requisição pela rede. Essa abordagem
elimina o \emph{overhead} de manutenção em períodos de inatividade,
sendo mais adequada para redes com padrões de tráfego esporádicos ou
imprevisíveis, características comuns em aplicações de sensoriamento. Os
protocolos híbridos, como o \emph{Zone Routing Protocol} (ZRP), combinam
características de ambas as abordagens: utilizam roteamento proativo
dentro de zonas locais e roteamento reativo para comunicação entre
zonas. Essa estratégia busca balancear as vantagens de ambas as
abordagens, mas introduz complexidade adicional na implementação. O Quadro~\ref{quad:protocolos-comparativo} sintetiza o comparativo entre as três categorias.

O Quadro~\ref{quadro:comp-protocolos} apresenta um comparativo entre as categorias de protocolos de roteamento ad-hoc.

{\def\LTcaptype{quadro}
\begin{longtable}[]{@{}
  >{\raggedright\arraybackslash}p{(\linewidth - 6\tabcolsep) * \real{0.2500}}
  >{\raggedright\arraybackslash}p{(\linewidth - 6\tabcolsep) * \real{0.2500}}
  >{\raggedright\arraybackslash}p{(\linewidth - 6\tabcolsep) * \real{0.2500}}
  >{\raggedright\arraybackslash}p{(\linewidth - 6\tabcolsep) * \real{0.2500}}@{}}
\caption{Comparativo entre categorias de protocolos de roteamento ad-hoc}\label{quadro:comp-protocolos}\\
\toprule\noalign{}
\begin{minipage}[b]{\linewidth}\raggedright
Característica
\end{minipage} & \begin{minipage}[b]{\linewidth}\raggedright
Proativo
\end{minipage} & \begin{minipage}[b]{\linewidth}\raggedright
Reativo
\end{minipage} & \begin{minipage}[b]{\linewidth}\raggedright
Híbrido
\end{minipage} \\
\midrule\noalign{}
\endhead
\bottomrule\noalign{}
\endlastfoot
Exemplos & OLSR, DSDV & AODV, DSR & ZRP \\
Descoberta de rotas & Antecipada (sempre) & Sob demanda & Combinada \\
Latência inicial & Baixa & Alta (descoberta) & Média \\
Overhead de controle & Constante (alto) & Variável (baixo) & Moderado \\
Uso de memória & Alto & Baixo/médio & Médio \\
Adequação para WSNs & Limitada & Boa & Moderada \\
\end{longtable}
}

Fonte: Adaptado de \citeonline{chaizeng2021} e \citeonline{priyadarshi2025}.

\subsection{Soluções existentes de mesh para
ESP32}\label{soluuxe7uxf5es-existentes-de-mesh-para-esp32}

O ecossistema ESP32 oferece algumas soluções de rede mesh, cada uma com
características distintas. O ESP-WIFI-MESH (também conhecido como
ESP-MDF) é a solução oficial da Espressif, implementando uma arquitetura
de árvore auto-organizada sobre a pilha Wi-Fi convencional. Embora
ofereça recursos avançados de gerenciamento de rede, o ESP-WIFI-MESH
impõe consumo energético elevado devido à utilização contínua do rádio
Wi-Fi para manutenção da topologia \cite{sestak2022}. A biblioteca
PainlessMesh, desenvolvida pela comunidade open-source, oferece uma
alternativa simplificada para criação de redes mesh com ESP32 e ESP8266.
Similar ao ESP-WIFI-MESH, a PainlessMesh opera sobre a pilha Wi-Fi,
herdando as mesmas limitações de consumo energético. Conforme observado
por \citeonline{sestak2022}, o uso do Wi-Fi padrão para manutenção da
topologia mesh resulta em \emph{overhead} significativo de processamento
e energia, comprometendo a aplicabilidade dessas soluções em
dispositivos alimentados por baterias. Abordagens alternativas têm
explorado a construção de redes mesh diretamente sobre o ESP-NOW.
\citeonline{cujilema2023} propuseram o algoritmo BRAM-NOW para criação
de redes mesh seguras, demonstrando latência média de 75 ms e taxa de
perda de pacotes de 9,25\% em ambiente residencial. Entretanto, essa
abordagem não se fundamenta em protocolos de roteamento padronizados,
limitando sua comparabilidade com outras soluções e dificultando
extensões futuras.

\section{O Protocolo AODV}\label{o-protocolo-aodv}

O \emph{Ad-hoc On-Demand Distance Vector} (AODV) é um protocolo de
roteamento reativo especificado na RFC 3561 \cite{perkins2003},
projetado para redes ad-hoc móveis (MANETs) com populações de dezenas a
milhares de nós. O protocolo oferece rápida adaptação a condições
dinâmicas de enlace, baixo \emph{overhead} de processamento e memória, e
utilização eficiente da rede, características que o tornam adequado para
adaptação em redes de sensores com recursos limitados.

\subsection{Mensagens de controle}\label{mensagens-de-controle}

O AODV define quatro tipos de mensagens de controle para descoberta e
manutenção de rotas: a) RREQ (\emph{Route Request}): mensagem de
requisição de rota, disseminada por \emph{broadcast} quando um nó
precisa descobrir um caminho para um destino desconhecido. Contém:
endereço IP de origem, número de sequência de origem, endereço IP de
destino, último número de sequência conhecido do destino, identificador
único da requisição (RREQ ID) e contador de saltos; b) RREP (\emph{Route
Reply}): mensagem de resposta, enviada em unicast de volta à origem
quando o destino ou um nó intermediário com rota válida recebe um RREQ.
Contém: endereço IP de destino, número de sequência de destino, endereço
IP de origem, contador de saltos e tempo de vida da rota; c) RERR
(\emph{Route Error}): mensagem de erro de rota, enviada quando uma
quebra de enlace é detectada. Notifica os nós afetados sobre destinos
que se tornaram inalcançáveis através do enlace interrompido; d)
RREP-ACK (\emph{Route Reply Acknowledgment}): mensagem de confirmação
opcional, utilizada para detectar enlaces unidirecionais quando
necessário.

\subsection{Processo de descoberta de
rotas}\label{processo-de-descoberta-de-rotas}

Quando um nó de origem necessita enviar dados para um destino para o
qual não possui rota válida em sua tabela de roteamento, inicia o
processo de descoberta de rota através da difusão de uma mensagem RREQ.
Antes da transmissão, o nó incrementa seu próprio número de sequência e
armazena o par (endereço de origem, RREQ ID) em um buffer temporal para
evitar reprocessamento de requisições duplicadas. Ao receber um RREQ,
cada nó intermediário: (i) cria ou atualiza uma entrada de rota reversa
para a origem; (ii) incrementa o contador de saltos; (iii) verifica se é
o destino ou se possui rota válida para o destino com número de
sequência adequado; (iv) caso positivo, gera e envia um RREP em unicast;
caso negativo, retransmite o RREQ para seus vizinhos (exceto aquele de
quem recebeu). A mensagem RREP é transmitida em unicast através do
caminho reverso estabelecido durante a propagação do RREQ. À medida que
o RREP percorre esse caminho, cada nó intermediário cria ou atualiza uma
entrada de rota direta para o destino, estabelecendo o próximo salto
(\emph{next hop}) como o vizinho de quem recebeu o RREP. Quando o RREP
alcança a origem, a rota bidirecional está estabelecida e a transmissão
de dados pode iniciar.

\subsection{Números de sequência e prevenção de
loops}\label{nuxfameros-de-sequuxeancia-e-prevenuxe7uxe3o-de-loops}

Uma característica distintiva do AODV é a utilização de números de
sequência de destino para garantir a ausência de loops de roteamento
(\emph{loop freedom}). Cada nó mantém seu próprio número de sequência,
incrementando-o antes de originar uma nova requisição de rota ou
resposta. Os números de sequência permitem distinguir informações de
roteamento atualizadas de informações obsoletas: entre duas rotas para
um mesmo destino, é sempre selecionada aquela com maior número de
sequência; em caso de empate, seleciona-se a rota com menor número de
saltos \cite{perkins2003}.

\subsection{Manutenção de rotas e tratamento de
falhas}\label{manutenuxe7uxe3o-de-rotas-e-tratamento-de-falhas}

O AODV monitora a conectividade com os próximos saltos das rotas ativas
através de dois mecanismos complementares: (i) mensagens HELLO
periódicas, que são RREPs com TTL=1 enviadas a cada HELLO\_INTERVAL
(tipicamente 1 segundo); e (ii) detecção de falha de enlace na camada de
enlace, quando disponível (por exemplo, ausência de ACK em transmissões
IEEE 802.11). Quando uma quebra de enlace é detectada, o nó que
identifica a falha invalida todas as rotas que utilizam o enlace
interrompido e gera uma mensagem RERR listando os destinos afetados. A
mensagem RERR é propagada para os precursores --- nós que utilizam o
remetente como próximo salto para os destinos afetados --- permitindo
que invalidem suas respectivas entradas de rota. Opcionalmente, o nó que
detecta a falha pode tentar reparo local da rota antes de propagar o
RERR.

\subsection{Parâmetros de
configuração}\label{paruxe2metros-de-configurauxe7uxe3o}

A RFC 3561 define valores padrão para diversos parâmetros configuráveis
que influenciam o comportamento e desempenho do protocolo. O Quadro~\ref{quadro:aodv-params}
apresenta os principais parâmetros e seus valores padrão, que podem ser
ajustados conforme as características específicas da rede.

{\def\LTcaptype{quadro}
\begin{longtable}[]{@{}lll@{}}
\caption{Parâmetros de configuração do AODV (RFC 3561)}\label{quadro:aodv-params}\\
\toprule\noalign{}
Parâmetro & Valor Padrão & Descrição \\
\midrule\noalign{}
\endhead
\bottomrule\noalign{}
\endlastfoot
ACTIVE\_ROUTE\_TIMEOUT & 3.000 ms & Tempo de vida de rota ativa \\
HELLO\_INTERVAL & 1.000 ms & Intervalo entre mensagens HELLO \\
ALLOWED\_HELLO\_LOSS & 2 & Perdas de HELLO antes de assumir falha \\
NET\_DIAMETER & 35 & Número máximo de saltos na rede \\
NODE\_TRAVERSAL\_TIME & 40 ms & Tempo estimado de travessia por nó \\
RREQ\_RETRIES & 2 & Tentativas de descoberta de rota \\
RREQ\_RATELIMIT & 10 & Máximo de RREQs por segundo \\
TTL\_START & 1 & TTL inicial para expanding ring search \\
TTL\_INCREMENT & 2 & Incremento de TTL a cada tentativa \\
TTL\_THRESHOLD & 7 & TTL máximo antes de usar NET\_DIAMETER \\
\end{longtable}
}

Fonte: RFC 3561 \cite{perkins2003}.
